% -*- root: ../presentation.tex -*-
\section{Results}

%==============================================================================
\begin{frame}
    \frametitle{Results: Cumulative oil production}
    \begin{itemize}
        \item Overall relatively good fit with results in paper.
        \pause
        \item Results for cumulative oil recovery after 1500:
        \begin{itemize}
             \item Results for cases 1-3 are generally lower than in the paper, with deviations ranging from $0.4-14.7\%$ with the average error being $6.45\%$. 
            \pause
             \item Numbers with the refined grid are lower than with the original grid. 
            \pause
             \item The deviations for the cases with short wells are, overall, significantly higher than the deviations for long wells.
            \pause
             \item Results for case 4 fit very well with the results in the paper, with deviations ranging from $0.4\%$ to $3.3\%$.
         \end{itemize} 
        \pause
        \item In the paper refined grid was not used for case 4. Some results for case 4 with the refined grid are shown here anyway.
    \end{itemize}
\end{frame}
%==============================================================================


%==============================================================================
\begin{frame}
    \frametitle{Results: Simulated Cumulative Oil Production}
    \footnotesize
    \centerline{% -*- root: ../../../plots.tex -*-

\begin{tikzpicture}
\pgfplotsset{width=1.1\textwidth, height=0.8\textheight}
\begin{axis}[
    ybar,
    enlargelimits=0.15,
    legend style={at={(0.5,-0.15)},
      anchor=north,legend columns=-1},
    ylabel={MSTB after 1500 days},
    symbolic x coords={1a,1b,2a,2b,3a,3b,4a,4b},
    xtick=data,
    bar width=7pt,
    ]
    \addplot coordinates {(1a,757) (1b,951) (2a,1034) (2b,1251) (3a,1229) (3b,1445) (4a,696) (4b,827) };
    \addplot coordinates {(1a,736) (1b,947) (2a,971)  (2b,1223) (3a,1099) (3b,1365) (4a,705) (4b,824) };
    \addplot coordinates {(1a,703) (1b,925) (2a,923)  (2b,1194) (3a,1049) (3b,1344) (4a,705) (4b,824) };
    \legend{Paper,Original grid,Refined grid}
\end{axis}
\end{tikzpicture}
}
\end{frame}
%==============================================================================


%==============================================================================
\begin{frame}
    \frametitle{Results: Simulated Cumulative Oil Production}
    \footnotesize
    \centerline{% -*- root: ../../../plots.tex -*-

\begin{tikzpicture}
\pgfplotsset{width=1.1\textwidth, height=0.8\textheight}
\begin{axis}[
    ybar,
    enlargelimits=0.15,
    legend style={at={(0.5,-0.15)},
      anchor=north,legend columns=-1},
    ylabel={\% Dev. in Cum. Oil Prod. after 1500 days},
    symbolic x coords={1a,1b,2a,2b,3a,3b,4a,4b},
    xtick=data,
    bar width=12pt,
    nodes near coords,
    nodes near coords align={vertical},
    ]
    \addplot coordinates {(1a,2.8) (1b,0.4) (2a,6.4)  (2b,3.2) (3a,11.9) (3b,5.9) (4a,3.3) (4b,0.4) };
    \addplot coordinates {(1a,7.2) (1b,2.8) (2a,10.6) (2b,4.5) (3a,14.7) (3b,7.0) (4a,1.2) (4b,0.4) };
    \legend{Original grid,Refined grid}
\end{axis}
\end{tikzpicture}
}
\end{frame}
%==============================================================================


%==============================================================================
\begin{frame}[label=oilprod13]
    \frametitle{Results: Oil production - Cases 1-3}
    \centerline{% -*- root: ../../../plots.tex -*-
\begin{tikzpicture}[font = \sffamily]
\pgfplotsset{legend style={legend columns=6}}
\pgfplotsset{legend style={at={(0.5,-0.30)}, font=\footnotesize}}
\pgfplotsset{width=1.0\textwidth, height=0.7\textheight}
    \begin{axis}[
        xlabel={Time (days)},
        ylabel={Cumulative Oil (stb)},
        xtick={200,400,...,1600},
        ymin=0, ymax=1.5e6,
        xmin=0, xmax=1500,
        axis y line*=left,
        xlabel near ticks,
        ylabel near ticks
        ]
        \addplot[color=red, mark=o] table[x=TIME, y=WOPT:PROD, col sep=comma, skip coords between index={0}{1}]{data/spe71a.csv};
        \addplot[color=red, mark=x] table[x=TIME, y=WOPT:PROD, col sep=comma, skip coords between index={0}{1}]{data/spe71b.csv};
        \addplot[color=green, mark=o] table[x=TIME, y=WOPT:PROD, col sep=comma, skip coords between index={0}{1}]{data/spe72a.csv};
        \addplot[color=green, mark=x] table[x=TIME, y=WOPT:PROD, col sep=comma, skip coords between index={0}{1}]{data/spe72b.csv};
        \addplot[color=blue, mark=o] table[x=TIME, y=WOPT:PROD, col sep=comma, skip coords between index={0}{1}]{data/spe73a.csv};
        \addplot[color=blue, mark=x] table[x=TIME, y=WOPT:PROD, col sep=comma, skip coords between index={0}{1}]{data/spe73b.csv};
    \end{axis}
    \begin{axis}[
        ymin=0, ymax=4000,
        xmin=0, xmax=1500,
        ylabel={Oil rate (stb/day)},
        hide x axis,
        axis y line*=right,
        ylabel near ticks,
        legend entries={1a,
                        1b,
                        2a,
                        2b,
                        3a,
                        3b},
        ]
        \addplot[color=red, mark=o] table[x=TIME, y=WOPR:PROD, col sep=comma, skip coords between index={0}{1}]{data/spe71a.csv};
        \addplot[color=red, mark=x] table[x=TIME, y=WOPR:PROD, col sep=comma, skip coords between index={0}{1}]{data/spe71b.csv};
        \addplot[color=green, mark=o] table[x=TIME, y=WOPR:PROD, col sep=comma, skip coords between index={0}{1}]{data/spe72a.csv};
        \addplot[color=green, mark=x] table[x=TIME, y=WOPR:PROD, col sep=comma, skip coords between index={0}{1}]{data/spe72b.csv};
        \addplot[color=blue, mark=o] table[x=TIME, y=WOPR:PROD, col sep=comma, skip coords between index={0}{1}]{data/spe73a.csv};
        \addplot[color=blue, mark=x] table[x=TIME, y=WOPR:PROD, col sep=comma, skip coords between index={0}{1}]{data/spe73b.csv};
        \addlegendimage{mark=o, red}
        \addlegendimage{mark=x, red}
        \addlegendimage{mark=o, green}
        \addlegendimage{mark=x, green}
        \addlegendimage{mark=o, blue}
        \addlegendimage{mark=x, blue}
      \end{axis}
\end{tikzpicture}
}
\end{frame}
%==============================================================================

%==============================================================================
\begin{frame}
    \frametitle{Results: Oil production - Cases 1-3}
    \centerline{% -*- root: ../../../plots.tex -*-
\begin{tikzpicture}[font = \sffamily]
\pgfplotsset{legend style={legend columns=6}}
\pgfplotsset{legend style={at={(0.5,-0.30)}, font=\footnotesize}}
\pgfplotsset{width=1.0\textwidth, height=0.85\textheight}
    \begin{axis}[
        xlabel={Time (days)},
        ylabel={Cumulative Oil (stb)},
        xtick={200,400,...,1600},
        ymin=0, 
        xmin=0, xmax=1500,
        axis y line*=left,
        xlabel near ticks,
        ylabel near ticks
        ]
        \addplot[color=red, mark=o] table[x=TIME, y=WOPT:PROD, col sep=comma, skip coords between index={0}{1}]{data/spe71a.csv};
        \addplot[color=red, mark=x] table[x=TIME, y=WOPT:PROD, col sep=comma, skip coords between index={0}{1}]{data/spe71b.csv};
        \addplot[color=green, mark=o] table[x=TIME, y=WOPT:PROD, col sep=comma, skip coords between index={0}{1}]{data/spe72a.csv};
        \addplot[color=green, mark=x] table[x=TIME, y=WOPT:PROD, col sep=comma, skip coords between index={0}{1}]{data/spe72b.csv};
        \addplot[color=blue, mark=o] table[x=TIME, y=WOPT:PROD, col sep=comma, skip coords between index={0}{1}]{data/spe73a.csv};
        \addplot[color=blue, mark=x] table[x=TIME, y=WOPT:PROD, col sep=comma, skip coords between index={0}{1}]{data/spe73b.csv};

        \addplot[color=red, mark=o, dashed] table[x=TIME, y=WOPT:PROD, col sep=comma, skip coords between index={0}{1}]{data/spe71a-refined.csv};
        \addplot[color=red, mark=x, dashed] table[x=TIME, y=WOPT:PROD, col sep=comma, skip coords between index={0}{1}]{data/spe71b-refined.csv};
        \addplot[color=green, mark=o, dashed] table[x=TIME, y=WOPT:PROD, col sep=comma, skip coords between index={0}{1}]{data/spe72a-refined.csv};
        \addplot[color=green, mark=x, dashed] table[x=TIME, y=WOPT:PROD, col sep=comma, skip coords between index={0}{1}]{data/spe72b-refined.csv};
        \addplot[color=blue, mark=o, dashed] table[x=TIME, y=WOPT:PROD, col sep=comma, skip coords between index={0}{1}]{data/spe73a-refined.csv};
        \addplot[color=blue, mark=x, dashed] table[x=TIME, y=WOPT:PROD, col sep=comma, skip coords between index={0}{1}]{data/spe73b-refined.csv};
    \end{axis}
    \begin{axis}[
        ymin=0, ymax=4000,
        xmin=0, xmax=1500,
        ylabel={Oil rate (stb/day)},
        hide x axis,
        axis y line*=right,
        ylabel near ticks,
        ]
        \addplot[color=red, mark=o] table[x=TIME, y=WOPR:PROD, col sep=comma, skip coords between index={0}{1}]{data/spe71a.csv};
        \addplot[color=red, mark=x] table[x=TIME, y=WOPR:PROD, col sep=comma, skip coords between index={0}{1}]{data/spe71b.csv};
        \addplot[color=green, mark=o] table[x=TIME, y=WOPR:PROD, col sep=comma, skip coords between index={0}{1}]{data/spe72a.csv};
        \addplot[color=green, mark=x] table[x=TIME, y=WOPR:PROD, col sep=comma, skip coords between index={0}{1}]{data/spe72b.csv};
        \addplot[color=blue, mark=o] table[x=TIME, y=WOPR:PROD, col sep=comma, skip coords between index={0}{1}]{data/spe73a.csv};
        \addplot[color=blue, mark=x] table[x=TIME, y=WOPR:PROD, col sep=comma, skip coords between index={0}{1}]{data/spe73b.csv};

        \addplot[color=red, mark=o, dashed] table[x=TIME, y=WOPR:PROD, col sep=comma, skip coords between index={0}{1}]{data/spe71a-refined.csv};
        \addplot[color=red, mark=x, dashed] table[x=TIME, y=WOPR:PROD, col sep=comma, skip coords between index={0}{1}]{data/spe71b-refined.csv};
        \addplot[color=green, mark=o, dashed] table[x=TIME, y=WOPR:PROD, col sep=comma, skip coords between index={0}{1}]{data/spe72a-refined.csv};
        \addplot[color=green, mark=x, dashed] table[x=TIME, y=WOPR:PROD, col sep=comma, skip coords between index={0}{1}]{data/spe72b-refined.csv};
        \addplot[color=blue, mark=o, dashed] table[x=TIME, y=WOPR:PROD, col sep=comma, skip coords between index={0}{1}]{data/spe73a-refined.csv};
        \addplot[color=blue, mark=x, dashed] table[x=TIME, y=WOPR:PROD, col sep=comma, skip coords between index={0}{1}]{data/spe73b-refined.csv};

      \end{axis}
\end{tikzpicture}
}
\end{frame}
%==============================================================================

%==============================================================================
\begin{frame}[label=wor13]
    \frametitle{Results: Water-Oil Ratio - Cases 1-3}
    \centerline{% -*- root: ../../../plots.tex -*-
\begin{tikzpicture}
\pgfplotsset{legend style={legend columns=6}}
\pgfplotsset{legend style={at={(0.5,-0.30)}, font=\footnotesize}}
\pgfplotsset{width=1.0\textwidth, height=0.7\textheight}
    \begin{axis}[
        xlabel={Time (days)},
        xtick={200,400,...,1600},
        ylabel={WOR (stb/stb)},
        xtick={200,400,...,1600},
        ymin=0,
        xmin=0, xmax=1500,
        ]        
        \addplot[color=red, mark=o] 
                table[  x=TIME, 
                        y expr={\thisrow{WWPR:PROD} / \thisrow{WOPR:PROD}}, 
                        col sep=comma, skip coords between index={0}{1}
                     ]
                {data/spe71a.csv}; 
        \addplot[color=red, mark=x] 
                table[  x=TIME, 
                        y expr={\thisrow{WWPR:PROD} / \thisrow{WOPR:PROD}}, 
                        col sep=comma, skip coords between index={0}{1}
                     ]
                {data/spe71b.csv};
        \addplot[color=green, mark=o] 
                table[  x=TIME, 
                        y expr={\thisrow{WWPR:PROD} / \thisrow{WOPR:PROD}}, 
                        col sep=comma, skip coords between index={0}{1}
                     ]
                {data/spe72a.csv};
        \addplot[color=green, mark=x] 
                table[  x=TIME, 
                        y expr={\thisrow{WWPR:PROD} / \thisrow{WOPR:PROD}}, 
                        col sep=comma, skip coords between index={0}{1}
                     ]
                {data/spe72b.csv};
        \addplot[color=blue, mark=o] 
                table[  x=TIME, 
                        y expr={\thisrow{WWPR:PROD} / \thisrow{WOPR:PROD}}, 
                        col sep=comma, skip coords between index={0}{1}
                     ]
                {data/spe73a.csv};
        \addplot[color=blue, mark=x] 
                table[  x=TIME, 
                        y expr={\thisrow{WWPR:PROD} / \thisrow{WOPR:PROD}}, 
                        col sep=comma, skip coords between index={0}{1}
                     ]
                {data/spe73b.csv};
        \addlegendentry{1a};
        \addlegendentry{1b};
        \addlegendentry{2a};
        \addlegendentry{2b};
        \addlegendentry{3a};
        \addlegendentry{3b};
    \end{axis}
\end{tikzpicture}
}
\end{frame}
%==============================================================================

%==============================================================================
\begin{frame}
    \frametitle{Results: Water-Oil Ratio - Cases 1-3}
    \centerline{% -*- root: ../../../plots.tex -*-
\begin{tikzpicture}
\pgfplotsset{legend style={legend columns=6}}
\pgfplotsset{legend style={at={(0.5,-0.30)}, font=\footnotesize}}
\pgfplotsset{width=1.0\textwidth, height=0.7\textheight}
    \begin{axis}[
        xlabel={Time (days)},
        xtick={200,400,...,1600},
        ylabel={WOR (stb/stb)},
        xtick={200,400,...,1600},
        ymin=0,
        xmin=0, xmax=1500,
        ]        
        \addplot[color=red, mark=o] 
                table[  x=TIME, 
                        y expr={\thisrow{WWPR:PROD} / \thisrow{WOPR:PROD}}, 
                        col sep=comma, skip coords between index={0}{1}
                     ]
                {data/spe71a.csv}; 
        \addplot[color=red, mark=x] 
                table[  x=TIME, 
                        y expr={\thisrow{WWPR:PROD} / \thisrow{WOPR:PROD}}, 
                        col sep=comma, skip coords between index={0}{1}
                     ]
                {data/spe71b.csv};
        \addplot[color=green, mark=o] 
                table[  x=TIME, 
                        y expr={\thisrow{WWPR:PROD} / \thisrow{WOPR:PROD}}, 
                        col sep=comma, skip coords between index={0}{1}
                     ]
                {data/spe72a.csv};
        \addplot[color=green, mark=x] 
                table[  x=TIME, 
                        y expr={\thisrow{WWPR:PROD} / \thisrow{WOPR:PROD}}, 
                        col sep=comma, skip coords between index={0}{1}
                     ]
                {data/spe72b.csv};
        \addplot[color=blue, mark=o] 
                table[  x=TIME, 
                        y expr={\thisrow{WWPR:PROD} / \thisrow{WOPR:PROD}}, 
                        col sep=comma, skip coords between index={0}{1}
                     ]
                {data/spe73a.csv};
        \addplot[color=blue, mark=x] 
                table[  x=TIME, 
                        y expr={\thisrow{WWPR:PROD} / \thisrow{WOPR:PROD}}, 
                        col sep=comma, skip coords between index={0}{1}
                     ]
                {data/spe73b.csv};


        \addplot[color=red, mark=o, dashed] 
                table[  x=TIME, 
                        y expr={\thisrow{WWPR:PROD} / \thisrow{WOPR:PROD}}, 
                        col sep=comma, skip coords between index={0}{1}
                     ]
                {data/spe71a-refined.csv}; 
        \addplot[color=red, mark=x, dashed] 
                table[  x=TIME, 
                        y expr={\thisrow{WWPR:PROD} / \thisrow{WOPR:PROD}}, 
                        col sep=comma, skip coords between index={0}{1}
                     ]
                {data/spe71b-refined.csv};
        \addplot[color=green, mark=o, dashed] 
                table[  x=TIME, 
                        y expr={\thisrow{WWPR:PROD} / \thisrow{WOPR:PROD}}, 
                        col sep=comma, skip coords between index={0}{1}
                     ]
                {data/spe72a-refined.csv};
        \addplot[color=green, mark=x, dashed] 
                table[  x=TIME, 
                        y expr={\thisrow{WWPR:PROD} / \thisrow{WOPR:PROD}}, 
                        col sep=comma, skip coords between index={0}{1}
                     ]
                {data/spe72b-refined.csv};
        \addplot[color=blue, mark=o, dashed] 
                table[  x=TIME, 
                        y expr={\thisrow{WWPR:PROD} / \thisrow{WOPR:PROD}}, 
                        col sep=comma, skip coords between index={0}{1}
                     ]
                {data/spe73a-refined.csv};
        \addplot[color=blue, mark=x, dashed] 
                table[  x=TIME, 
                        y expr={\thisrow{WWPR:PROD} / \thisrow{WOPR:PROD}}, 
                        col sep=comma, skip coords between index={0}{1}
                     ]
                {data/spe73b-refined.csv};
    \end{axis}
\end{tikzpicture}
}
\end{frame}
%==============================================================================

%==============================================================================
\begin{frame}
    \frametitle{Results: Cum. Water Production - Cases 1-3}
    \centerline{\input{figures/plots/water/cumwat.tikz}}
\end{frame}
%==============================================================================

%==============================================================================
\begin{frame}
    \frametitle{Results: Cum. Water Production - Cases 1-3}
    \centerline{\input{figures/plots/water/cumwat-refined.tikz}}
\end{frame}
%==============================================================================

%==============================================================================
\begin{frame}
    \frametitle{Results: Coning}
    \begin{itemize}
        \item WOR ratio is higher for all cases with short producing wells.
        \item Likely due to a greater tendency for coning.
        \item The cumulative water production will still be greater for the long wells because they produce more fluid overall.
        \pause
        \item Results from refined grid show an overall higher WOR.
        \begin{itemize}
            \item Possibly more accurate simulation of coning tendencies.
            \item Contributes to the lower oil recovery.
        \end{itemize}
    \end{itemize}
\end{frame}
%============================================================================== 


%==============================================================================
\begin{frame}
    \frametitle{Results: Case 4}
    \begin{itemize}
        \item In case 4 less liquid is injected than is being produced.
        \pause
        \item This causes the pressure to drop, leading to gas being released from solution.
        \item Gas production drops to $\approx 0$ towards the end.
        \pause
        \item Oil and water prodction ``settles'' at a lower rate compared to case 3 due to lower reservoir pressure.
    \end{itemize}
    \pause
    \begin{remark}
        Focus here is on case 4b. Case 4a has the same trends, but, as with cases 1-3, somewhat lower OPR and higher WOR.
    \end{remark}
\end{frame}
%============================================================================== 


%==============================================================================
\begin{frame}
    \frametitle{Results: Oil production - Case 4}
    \centerline{% -*- root: ../../../plots.tex -*-
\begin{tikzpicture}[font = \sffamily]
\pgfplotsset{legend style={legend columns=2}}
\pgfplotsset{legend style={at={(0.5,-0.30)}, font=\footnotesize}}
\pgfplotsset{width=1.0\textwidth, height=0.7\textheight}
    \begin{axis}[
        xlabel={Time (days)},
        xtick={200,400,...,1600},
        ylabel={Cumulative Oil (stb)},
        ymin=0,
        xmin=0, xmax=1500,
        axis y line*=left,
        xlabel near ticks,
        ylabel near ticks
        ]
        \addplot[color=red, mark=o] table[x=TIME, y=WOPT:PROD, col sep=comma, skip coords between index={0}{1}]{data/spe74a.csv};
        \addplot[color=blue, mark=o] table[x=TIME, y=WOPT:PROD, col sep=comma, skip coords between index={0}{1}]{data/spe74b.csv};

        \addplot[color=red, mark=o, dashed] table[x=TIME, y=WOPT:PROD, col sep=comma, skip coords between index={0}{1}]{data/spe74a-refined.csv};
        \addplot[color=blue, mark=o, dashed] table[x=TIME, y=WOPT:PROD, col sep=comma, skip coords between index={0}{1}]{data/spe74b-refined.csv};
    \end{axis}
    \begin{axis}[
        ylabel={Oil rate (stb/day)},
        ymin=0,
        xmin=0, xmax=1500,
        hide x axis,
        axis y line*=right,
        ylabel near ticks,
        legend entries={$L=900$: Orig. grid,
                        $L=2100$: Orig. grid,
                        $L=900$: Ref. grid,
                        $L=2100$: Ref. grid},
        ]
        \addplot[color=red, mark=o] table[x=TIME, y=WOPR:PROD, col sep=comma, skip coords between index={0}{1}]{data/spe74a.csv};
        \addplot[color=blue, mark=o] table[x=TIME, y=WOPR:PROD, col sep=comma, skip coords between index={0}{1}]{data/spe74b.csv};

        \addplot[color=red, mark=o, dashed] table[x=TIME, y=WOPR:PROD, col sep=comma, skip coords between index={0}{1}]{data/spe74a-refined.csv};
        \addplot[color=blue, mark=o, dashed] table[x=TIME, y=WOPR:PROD, col sep=comma, skip coords between index={0}{1}]{data/spe74b-refined.csv};
      \end{axis}
\end{tikzpicture}
}
\end{frame}
%==============================================================================

% %==============================================================================
% \begin{frame}
%     \frametitle{Results: Water Production Rate - Case 4}
%     \centerline{\input{figures/plots/water/case4-wpr.tikz}}
% \end{frame}
% %==============================================================================

% %==============================================================================
% \begin{frame}
%     \frametitle{Results: Gas Production Rate - Case 4}
%     \centerline{\input{figures/plots/gas/case4-gpr.tikz}}
% \end{frame}
% %==============================================================================

%==============================================================================
\begin{frame}
    \frametitle{Results: Fluid Production Rates - Case 4b}
    \centerline{% -*- root: ../../../plots.tex -*-
\begin{tikzpicture}[font = \sffamily]
\pgfplotsset{legend style={legend columns=2}}
\pgfplotsset{legend style={at={(0.5,-0.30)}, font=\footnotesize}}
\pgfplotsset{width=1.0\textwidth, height=0.7\textheight}
    \begin{axis}[
        xlabel={Time (days)},
        xtick={200,400,...,1600},
        ylabel={Cumulative Oil (stb)},
        ymin=0,
        xmin=0, xmax=1500,
        axis y line*=left,
        xlabel near ticks,
        ylabel near ticks
        ]
        \addplot[color=red, mark=o] table[x=TIME, y=WOPT:PROD, col sep=comma, skip coords between index={0}{1}]{data/spe74a.csv};
        \addplot[color=blue, mark=o] table[x=TIME, y=WOPT:PROD, col sep=comma, skip coords between index={0}{1}]{data/spe74b.csv};

        \addplot[color=red, mark=o, dashed] table[x=TIME, y=WOPT:PROD, col sep=comma, skip coords between index={0}{1}]{data/spe74a-refined.csv};
        \addplot[color=blue, mark=o, dashed] table[x=TIME, y=WOPT:PROD, col sep=comma, skip coords between index={0}{1}]{data/spe74b-refined.csv};
    \end{axis}
    \begin{axis}[
        ylabel={Oil rate (stb/day)},
        ymin=0,
        xmin=0, xmax=1500,
        hide x axis,
        axis y line*=right,
        ylabel near ticks,
        legend entries={$L=900$: Orig. grid,
                        $L=2100$: Orig. grid,
                        $L=900$: Ref. grid,
                        $L=2100$: Ref. grid},
        ]
        \addplot[color=red, mark=o] table[x=TIME, y=WOPR:PROD, col sep=comma, skip coords between index={0}{1}]{data/spe74a.csv};
        \addplot[color=blue, mark=o] table[x=TIME, y=WOPR:PROD, col sep=comma, skip coords between index={0}{1}]{data/spe74b.csv};

        \addplot[color=red, mark=o, dashed] table[x=TIME, y=WOPR:PROD, col sep=comma, skip coords between index={0}{1}]{data/spe74a-refined.csv};
        \addplot[color=blue, mark=o, dashed] table[x=TIME, y=WOPR:PROD, col sep=comma, skip coords between index={0}{1}]{data/spe74b-refined.csv};
      \end{axis}
\end{tikzpicture}
}
\end{frame}
%==============================================================================

%==============================================================================
\begin{frame}
    \frametitle{Results: Pressure - All cases}
    \centerline{\input{figures/plots/pressure/combined-bhp.tikz}}
\end{frame}
%==============================================================================

%==============================================================================
\begin{frame}
    \frametitle{Results: Production profile development}
    \begin{itemize}
        \item The gas production rate drops until $\approx 200$ days. At this point $P_{bp}$ seems to have been reached and the gas production rate increases dramatically.
        \pause
        \item After $\approx 600$ days, at peak gas production, all gas seems to have gone out of solution.
        \begin{itemize}
            \item Pressure settles at the set minimum BHP (1500 psia).
            \item Gas production rate drops very quickly as the remaining free gas is produced.
            \item Water production rate drops equally fast due to low reservoir pressure and the fact that there is very little flow to facilitate coning behavior.
            \item At the same time oil production drops to $\approx 0$.
        \end{itemize}
        \pause
        \item Subsequently the water and oil production rates rise.
    \end{itemize}
\end{frame}
%==============================================================================

% %==============================================================================
% \begin{frame}
%     \frametitle{Results: Gas-Oil Ratio - Case 4}
%     \centerline{\input{figures/plots/gas/case4-gor.tikz}}
% \end{frame}
% %==============================================================================

% %==============================================================================
% \begin{frame}
%     \frametitle{Results: Gas-Oil Ratio - Cases 1-3}
%     \centerline{% -*- root: ../../../plots.tex -*-
\begin{tikzpicture}[font = \sffamily]
\pgfplotsset{legend style={legend columns=6}}
\pgfplotsset{legend style={at={(0.5,-0.30)}, font=\footnotesize}}
\pgfplotsset{width=1.0\textwidth, height=0.7\textheight}
    \begin{axis}[
        xlabel={Time (days)},
        ylabel={Cumulative Oil (stb)},
        xtick={200,400,...,1600},
        ymin=0, ymax=1.5e6,
        xmin=0, xmax=1500,
        axis y line*=left,
        xlabel near ticks,
        ylabel near ticks
        ]
        \addplot[color=red, mark=o] table[x=TIME, y=WOPT:PROD, col sep=comma, skip coords between index={0}{1}]{data/spe71a.csv};
        \addplot[color=red, mark=x] table[x=TIME, y=WOPT:PROD, col sep=comma, skip coords between index={0}{1}]{data/spe71b.csv};
        \addplot[color=green, mark=o] table[x=TIME, y=WOPT:PROD, col sep=comma, skip coords between index={0}{1}]{data/spe72a.csv};
        \addplot[color=green, mark=x] table[x=TIME, y=WOPT:PROD, col sep=comma, skip coords between index={0}{1}]{data/spe72b.csv};
        \addplot[color=blue, mark=o] table[x=TIME, y=WOPT:PROD, col sep=comma, skip coords between index={0}{1}]{data/spe73a.csv};
        \addplot[color=blue, mark=x] table[x=TIME, y=WOPT:PROD, col sep=comma, skip coords between index={0}{1}]{data/spe73b.csv};
    \end{axis}
    \begin{axis}[
        ymin=0, ymax=4000,
        xmin=0, xmax=1500,
        ylabel={Oil rate (stb/day)},
        hide x axis,
        axis y line*=right,
        ylabel near ticks,
        legend entries={1a,
                        1b,
                        2a,
                        2b,
                        3a,
                        3b},
        ]
        \addplot[color=red, mark=o] table[x=TIME, y=WOPR:PROD, col sep=comma, skip coords between index={0}{1}]{data/spe71a.csv};
        \addplot[color=red, mark=x] table[x=TIME, y=WOPR:PROD, col sep=comma, skip coords between index={0}{1}]{data/spe71b.csv};
        \addplot[color=green, mark=o] table[x=TIME, y=WOPR:PROD, col sep=comma, skip coords between index={0}{1}]{data/spe72a.csv};
        \addplot[color=green, mark=x] table[x=TIME, y=WOPR:PROD, col sep=comma, skip coords between index={0}{1}]{data/spe72b.csv};
        \addplot[color=blue, mark=o] table[x=TIME, y=WOPR:PROD, col sep=comma, skip coords between index={0}{1}]{data/spe73a.csv};
        \addplot[color=blue, mark=x] table[x=TIME, y=WOPR:PROD, col sep=comma, skip coords between index={0}{1}]{data/spe73b.csv};
        \addlegendimage{mark=o, red}
        \addlegendimage{mark=x, red}
        \addlegendimage{mark=o, green}
        \addlegendimage{mark=x, green}
        \addlegendimage{mark=o, blue}
        \addlegendimage{mark=x, blue}
      \end{axis}
\end{tikzpicture}
}
% \end{frame}
% %==============================================================================

%==============================================================================
\begin{frame}
    \frametitle{Results: Cum. Gas Prod. - All cases}
    \centerline{% -*- root: ../../../plots.tex -*-
\begin{tikzpicture}
\pgfplotsset{legend style={legend columns=5}}
\pgfplotsset{legend style={at={(0.5,-0.30)}, font=\footnotesize}}
\pgfplotsset{width=1.0\textwidth, height=0.7\textheight}
    \begin{axis}[
        xlabel={Time (days)},
        xtick={200,400,...,1600},
        ylabel={Cum. Gas Prod. (scf/day)},
        xtick={200,400,...,1600},
        ymin=0,
        xmin=0, xmax=1500,
        ]        
        \addplot[color=red, mark=o] table[  x=TIME, y=FGPT, col sep=comma, skip coords between index={0}{1}]{data/spe71a.csv}; 
        \addlegendentry{1a};
        
        \addplot[color=red, mark=x] table[  x=TIME, y=FGPT, col sep=comma, skip coords between index={0}{1}]{data/spe71b.csv}; 
        \addlegendentry{1b};

        \addplot[color=green, mark=o] table[  x=TIME, y=FGPT, col sep=comma, skip coords between index={0}{1}]{data/spe72a.csv}; 
        \addlegendentry{2a};

        \addplot[color=green, mark=x] table[  x=TIME, y=FGPT, col sep=comma, skip coords between index={0}{1}]{data/spe72b.csv}; 
        \addlegendentry{2b};

        \addplot[color=blue, mark=o] table[  x=TIME, y=FGPT, col sep=comma, skip coords between index={0}{1}]{data/spe73a.csv}; 
        \addlegendentry{3a};

        \addplot[color=blue, mark=x] table[  x=TIME, y=FGPT, col sep=comma, skip coords between index={0}{1}]{data/spe73b.csv}; 
        \addlegendentry{3b};

        \addplot[color=gray, mark=o] table[  x=TIME, y=FGPT, col sep=comma, skip coords between index={0}{1}]{data/spe74a.csv}; 
        \addlegendentry{4a};

        \addplot[color=gray, mark=x] table[  x=TIME, y=FGPT, col sep=comma, skip coords between index={0}{1}]{data/spe74b.csv}; 
        \addlegendentry{4b};

        \addplot[color=gray, mark=o, dashed] table[  x=TIME, y=FGPT, col sep=comma, skip coords between index={0}{1}]{data/spe74a-refined.csv}; 
        \addlegendentry{4a (r)};

        \addplot[color=gray, mark=x, dashed] table[  x=TIME, y=FGPT, col sep=comma, skip coords between index={0}{1}]{data/spe74b-refined.csv}; 
        \addlegendentry{4b (r)};
    \end{axis}
\end{tikzpicture}
}
\end{frame}
%==============================================================================

%==============================================================================
\begin{frame}
    \frametitle{Results: Pressure Drop Along Wellbore - Case 4}
    \centerline{\input{figures/plots/pressure/case4-drop.tikz}}
\end{frame}
%==============================================================================

%==============================================================================
\begin{frame}
    \frametitle{Results: Pressure Drop Along Wellbore - Case 4}
    \begin{itemize}
        \item The highest pressure drop coincides with the time of peak gas production.
        \item The very high flow rate ($\approx 3.2 \times 10^4$ scf/day) causes high friction.
    \end{itemize}
\end{frame}
%==============================================================================

%==============================================================================
\begin{frame}
    \frametitle{Results: Pressure Drop Along Wellbore - Case 4}
    Possible explanation for higher pressure difference in the refined model:
    \begin{itemize}
        \item Refined model has more blocks on the ends of the well.
        \begin{itemize}
            \item The end blocks are further from each other in the refined model than in the original model.
            \item The pressure in the end blocks in the original model corresponds to the average of the pressure in the \emph{two} end blocks in the refined model.
            \item In the case of pressure difference between heel and toe, this is the same as having a (slightly) longer well.
        \end{itemize}
    \end{itemize}
\end{frame}
%==============================================================================

% %==============================================================================
% \begin{frame}
%     \frametitle{Results: Fluid Production Rates - Case 3b \& 4b}
%     \centerline{% -*- root: ../../../plots.tex -*-
\begin{tikzpicture}[font = \sffamily]
\pgfplotsset{legend style={legend columns=3}}
\pgfplotsset{legend style={at={(0.5,-0.30)}, font=\footnotesize}}
\pgfplotsset{width=1.0\textwidth, height=0.7\textheight}
    \begin{axis}[
        xlabel={Time (days)},
        xtick={200,400,...,1600},
        ylabel={Gas Prod. Rate (scf/day)},
        ymin=0,
        xmin=0, xmax=1500,
        axis y line*=left,
        xlabel near ticks,
        ylabel near ticks
        ]
        \addplot[color=red, mark=o] table[x=TIME, y=FGPR, col sep=comma, skip coords between index={0}{1}]{data/spe74b.csv};
        \addplot[color=red, mark=o, dashed] table[x=TIME, y=FGPR, col sep=comma, skip coords between index={0}{1}]{data/spe73b.csv};
    \end{axis}
    \begin{axis}[
        ylabel={Oil/Water Prod. Rate (stb/day)},
        ymin=0,
        xmin=0, xmax=1500,
        hide x axis,
        axis y line*=right,
        ylabel near ticks,
        legend entries={Gas,
                        Oil,
                        Water},
        ]
        \addplot[color=red, mark=o] table[x=TIME, y=FOPR, col sep=comma, skip coords between index={0}{1}]{data/spe74b.csv};
        \addplot[color=green, mark=o] table[x=TIME, y=FOPR, col sep=comma, skip coords between index={0}{1}]{data/spe74b.csv};
        \addplot[color=blue, mark=o] table[x=TIME, y=FWPR, col sep=comma, skip coords between index={0}{1}]{data/spe74b.csv};

        \addplot[color=red, mark=o, dashed] table[x=TIME, y=FOPR, col sep=comma, skip coords between index={0}{1}]{data/spe73b.csv};
        \addplot[color=green, mark=o, dashed] table[x=TIME, y=FOPR, col sep=comma, skip coords between index={0}{1}]{data/spe73b.csv};
        \addplot[color=blue, mark=o, dashed] table[x=TIME, y=FWPR, col sep=comma, skip coords between index={0}{1}]{data/spe73b.csv};
      \end{axis}
\end{tikzpicture}
}
% \end{frame}
% %==============================================================================

%==============================================================================
\begin{frame}
    \frametitle{Results: Water-Oil Ratio - Case 4}
    \centerline{% -*- root: ../../../plots.tex -*-
\begin{tikzpicture}
\pgfplotsset{legend style={legend columns=2}}
\pgfplotsset{legend style={at={(0.5,-0.30)}, font=\footnotesize}}
\pgfplotsset{width=1.0\textwidth, height=0.7\textheight}
    \begin{axis}[
        xlabel={Time (days)},
        xtick={200,400,...,1600},
        ylabel={WOR (stb/stb)},
        xtick={200,400,...,1600},
        ymin=0,
        xmin=0, xmax=1500,
        ]        
        \addplot[color=red, mark=o] 
                table[  x=TIME, 
                        y expr={\thisrow{WWPR:PROD} / \thisrow{WOPR:PROD}}, 
                        col sep=comma, skip coords between index={0}{1}
                     ]
                {data/spe74a.csv}; 
        \addplot[color=red, mark=o, dashed] 
                table[  x=TIME, 
                        y expr={\thisrow{WWPR:PROD} / \thisrow{WOPR:PROD}}, 
                        col sep=comma, skip coords between index={0}{1}
                     ]
                {data/spe74a-refined.csv};

        \addplot[color=blue, mark=x] 
                table[  x=TIME, 
                        y expr={\thisrow{WWPR:PROD} / \thisrow{WOPR:PROD}}, 
                        col sep=comma, skip coords between index={0}{1}
                     ]
                {data/spe74b.csv};
 
        \addplot[color=blue, mark=x, dashed] 
                table[  x=TIME, 
                        y expr={\thisrow{WWPR:PROD} / \thisrow{WOPR:PROD}}, 
                        col sep=comma, skip coords between index={0}{1}
                     ]
                {data/spe74b-refined.csv};

        \addlegendentry{4a: Original grid};
        \addlegendentry{4a: Refined grid};
        \addlegendentry{4b: Original grid};
        \addlegendentry{4b: Refined grid};
    \end{axis}
\end{tikzpicture}
}
\end{frame}
%==============================================================================

%==============================================================================
\begin{frame}
    \frametitle{Results: Cum. Water Production - Case 4}
    \centerline{% -*- root: ../../../plots.tex -*-
\begin{tikzpicture}
\pgfplotsset{legend style={legend columns=4}}
\pgfplotsset{legend style={at={(0.5,-0.30)}, font=\footnotesize}}
\pgfplotsset{width=1.0\textwidth, height=0.7\textheight}
    \begin{axis}[
        xlabel={Time (days)},
        xtick={200,400,...,1600},
        ylabel={Cum. Wat. Prod (stb/day)},
        xtick={200,400,...,1600},
        ymin=0,
        xmin=0, xmax=1500,
        ]        
        \addplot[color=gray, mark=o] table[  x=TIME, y=FOPT, col sep=comma, skip coords between index={0}{1}]{data/spe74a.csv}; 
        \addlegendentry{4a};

        \addplot[color=gray, mark=x] table[  x=TIME, y=FOPT, col sep=comma, skip coords between index={0}{1}]{data/spe74b.csv}; 
        \addlegendentry{4b};
    \end{axis}
\end{tikzpicture}
}
\end{frame}
%==============================================================================

%==============================================================================
\begin{frame}
    \frametitle{Results: Case 4 - Coning}
    \begin{itemize}
        \item Coning tendencies are less easy to spot from WOR.
        \item Short well has a higher WOR while the gas production is relatively low; after that things change.
    \end{itemize}
\end{frame}
%==============================================================================
